% Merge summary and introduction 
\section{Rubin Early Science Program}

Community expectations for early science with Rubin are high due to the transformative nature of the LSST data and the densely-sampled observations obtained during the commissioning period.
Rubin Observatory's \emph{Early Science Program} was established to provide the Rubin community with early access to the data products and services needed to produce high-impact science during the period from commissioning up to, but not including, Data Release 1 (DR1).

\subsection{Early Science Definitions}
\label{ssec:defn}

The following terms and their definitions are used in this document:
\begin{itemize}

\item\textbf{Commissioning:} The Rubin Construction Commissioning period, from LSSTComCam First Photon on 2024-10-24 through the end of Construction, \cstrcompdate.
\item\textbf{Early Operations:} The period from the start of Rubin Operations, \opsstartdate,  through the first Data Release, DR1.

\item\textbf{Rubin Operations Early Optimization Phase:} The period from the start of Rubin Operations, \opsstartdate, through the start of the LSST, \lsststartdate.

\item\textbf{Early Science:} Any science enabled by Rubin for its community during the period from the beginning of Commissioning through the first Data Release, DR1.
\end{itemize}

\subsection{Motivation for an Early Science Program} 
 \label{ssec:motivation}
The Early Science program is motivated by the desire to:
\begin{itemize}
\item enable high-impact science with Rubin Observatory data as early as possible;
\item provide early access to both static-sky and time-domain science-ready data products to support the community to prepare for science with the LSST;
\item enable early time-domain astronomy via early Alert Production; and 
\item help drive the development of Rubin operations capabilities prior to survey start and prepare the team to be operations-ready.
\end{itemize}

\subsection{Elements of the Early Science Program}
 \label{ssec:elements}
The Early Science Program consists of the following elements:
\begin{itemize}
	\item A series of three \textbf{Data Previews (DP)}, DP0, DP1 and DP2,  based on either simulated LSST-like data or data taken during the Rubin Observatory Commissioning period and Early Optimization Phase. 
	\item A world-public \textbf{Stream of Alerts} from transient, variable, and moving sources, which will be scaled up continuously during Commissioning and the Early Optimization Phase, and then continue throughout Operations.
	\item  \textbf{Template generation}, beginning before regular survey operations based on data collected during  Commissioning and the Early Optimization Phase  with LSSTCam,  and then continuing as needed until DRP templates are available for the entire sky.
The aim is to maximize the number of templates available for Alert Production in year 1 of the LSST.
	\item \textbf{Nightly Processed Visit Images (PVIs)} and associated direct-image \texttt{Source} catalogs of detections up until  Data Release 1. 
\end{itemize}

Before the Early Science Program was conceived, the project planned at least two data releases in the first year of operations.
The first of these, DR1, was to be based on the first six months of survey data and served as the original early preview of LSST data.
The Early Science Program, introduced later as a much more comprehensive plan for releasing early data, initially incorporated this six-month DR1.
However, its schedule overlapped with that of DP2, and two Data Release Processing campaigns cannot be in flight at the same time.
The six-month DR1 was therefore cancelled.
DR1 is now defined as the first data release based on a full year of LSST data.
Because it is a full data release rather than a preview, it falls outside the scope of the Early Science Program (\S\ref{ssec:dr1}).
DR1 will be led by the Rubin Operations Data Release Board, and its planning and contents will be documented in \citeds{RDO-011}.

%
\subsection{Transition to Operations and Early Science}
\label{ssec:transition}
The Rubin Construction Project was declared substantially complete on \cstrcompdate, marking the delivery of an integrated system capable of capturing, transferring, and processing science-grade images consistent with the Rubin/LSST Science Requirements Document (SRD; \citealt[][]{LPM-17}).
The following night, the observatory began on-sky operations and entered the Rubin Operations Early Optimization Phase.
Commissioning demonstrated that the as-built system could meet the LSST science requirements; from the start of Rubin Operations through the start of the LSST survey, \lsststartdate, the Operations team focused on achieving stable, repeatable nightly performance while coordinating remaining activities with the Construction team.

Data collected during Commissioning (\S~\ref{sec:commissioning}) and Early Operations provide the foundation for Early Science data products, offering the community an exceptional early dataset as the survey begins systematic sky coverage leading up to Data Release 1 (DR1).
While the Operations team aims to deliver the broadest feasible range of early data products and services, the contents and timelines for Data Previews, the early Alert Stream, and supporting services remain subject to change and reflect the best available estimates at the time of publication.

Rubin will endeavor provide at least 4 weeks advance notice before a data release or new data access service goes live.
For alerts, as this is an ongoing program, alerts will continue to ramp up following initial release, and no new announcements will be made for each subsequent increase.
Given the rapidly evolving nature of the commissioning and Early Science periods, and the associated uncertainties, it may not always be possible to provide four weeks' notice; however, we will endeavor to do so whenever possible.


\subsection{Factors Impacting the Early Science Program}
\label{ssec:impact}
Factors affecting the schedule and contents of the Early Science program can be broadly grouped into technical considerations and policy considerations.
Technical considerations include:
\begin{itemize}
\item The operational status of the observatory during Commissioning  and the Early Optimization Phase;
\item The nature and quality of the data collected during Commissioning and the Early Optimization Phase;
\item The readiness of the data processing pipelines, and data distribution and access services.
\end{itemize}
Policy factors include:
\begin{itemize}
\item The 30-day embargo on all pixel data during Commissioning (concluded);
\item The 80-hour embargo on all pixel data throughout the full duration of Operations and the LSST (current);
\item The  construction security review, which must be successfully completed prior to the release of any Prompt data products (successfully completed);
\item The Rubin First Look (RFL) media event \citep[][]{RTN-083},  before which no Rubin image data could be released (concluded).
\end{itemize}
Of the above policy factors, only the 80-hour embargo on all pixel data still remains active, and will continue to be, for the duration of Rubin Operations.

In this document, the term ``stretch goal'' will be used to describe cases where any uncertainty is due to a technical or scientific consideration and ``TBD'' (To Be Decided) will be used when the influencing factor is of a policy nature.